% SEGMENT OLP-0686-S01
% Part: proof-theory
% Chapter: propositions-as-types
% Section: introduction

\documentclass[../../../include/open-logic-section]{subfiles}

\begin{document}

\olfileid{int}{pty}{int}

\olsection{அறிமுகம்}

லாம்டா கலனம் கணக்கீட்டின் ஓர் எளிய மாதிரி. மாறிகளிலிருந்து
கட்டப்பட்ட சொற்களின்மீது அது செயல்படுகிறது. $N$ உம் $M$
உம் சொற்கள் என்றால், $(NM)$ உம் ஒரு சொல்; இது
``$N$ ஐ~$M$ மீது செயல்படுத்தல்'' எனப் பொருள்படும்.
$N$ ஒரு சொல் என்றால், $\lambd[x][N]$ உம் ஒரு சொல்.
$N$ இல் மாறி~$x$ இருந்தால், $\lambd[x][N]$ இல்
அந்தச் சொல்லிலிருந்து வந்த $x$ இன் கட்டற்ற நிகழ்வுகள்
$\lambd[x]$ செயற்குறியால் கட்டுப்படுகின்றன; ஆகவே அவை
இனி கட்டற்றவை அல்ல. $(\lambd[x][N]M)$ என்ற வடிவச் சொல்
$\Subst{N}{M}{x}$ ஆக $\beta$-சுருங்கும் என்போம்.
இங்கு $N$ இல் கட்டற்றதாக வரும் ஒவ்வொரு $x$ ஐயும்
$M$ ஆல் மாற்றியதே அந்த விளைவு.
$N$ இன் ஒரு உட்சொல்லை $\beta$-சுருக்குவதால் $N'$
கிடைத்தால், $N$ ஆனது~$N'$ ஆக ஒரு $\beta$-குறைத்தல்
படியில் மாறுகிறது; அதை $N \redone N'$ என எழுதுகிறோம்.
$\lambd$-கலனத்தில் ஒரு கணக்கீடு என்பது
$N \redone N_1 \redone N_2 \redone \dots$ என்ற தொடர்.
இந்தத் தொடர், எந்த உட்சொல்லும் $\beta$-சுருங்க முடியாத
$N_k$ என்ற சொல்லில் முடிந்தால், அதனை $N$ இன்
\emph{இயல்புவடிவம்} என்போம். $\lambd$-கலனக் கணக்கீட்டை
இவ்வகைத் தொடர்களாகப் புரிந்துகொள்ளலாம். இயல்புவடிவம்
கிடைத்தால் கணக்கீடு நிற்கிறது; அந்த இயல்புவடிவமே
கணக்கீட்டின் விளைவு. வகைக் கட்டுப்பாடு இல்லாத இந்த
எளிய கணக்கீட்டு மாதிரி டியூரிங்-முழுமையானது: ஒவ்வொரு
பகுதி மீள்வரையறைச் சார்பையும் $\lambd$-கலனத்தின்
ஒரு சொல்லால் கணக்கிடலாம்.

% SEGMENT OLP-0686-S02
ஹாஸ்கல் கரியும் வில்லியம் ஹாவர்டும், இயல்பான
வருவித்தலுக்கும் $\lambd$-கலனத்துக்கும் இடையிலான
நெருங்கிய ஒற்றுமையைத் தனித்தனியாகக் கண்டறிந்தனர்.
உள்ளுணர்வாக, $\lambd[x][N]$ என்ற சொல் ஒரு சார்பு:
$x$ அதன் உள்ளீடு; கொடுக்கப்பட்ட~$x$ க்குரிய மதிப்பை
எப்படிக் கணக்கிடுவது என்பதை $N$ தெரிவிக்கிறது.
$(\lambd[x][N]M)$ என்பது $\lambd[x][N]$ சார்பை
$M$ என்ற உள்ளீட்டின்மீது செயல்படுத்துவது. இதை எப்படி
மதிப்பிடுவது என்பதை $\beta$-சுருக்கல் சொல்கிறது:
$N$ இல் $x$ ஐ~$M$ ஆல் மாற்றி, கிடைத்த சொல்லைத்
தொடர்ந்து மதிப்பிட வேண்டும்.

இப்போது இச்சார்புகளின் உள்ளீட்டுத் தொகுதியும்
மதிப்புகள் அடங்கும் இலக்குத் தொகுதியும் நிலையானவை
எனக் கொள்வோம். $\lambd[x][N]$ இன் வரையறுக்களம் $A$
என்றும் அதன் மதிப்புகள் அடங்கும் இலக்குத் தொகுதி~$B$
என்றும் இருந்தால், மாறி~$x$ ஆனது~$A$ இலிருந்து
மட்டுமே உள்ளீடுகளைப் பெறும்; $N$ ஆனது~$B$ இல்
மதிப்புகளைத் தரும். எனவே $\lambd[x][N]$ ஒரு
$A \to B$ சார்பாகவும் $M \in A$ ஆகவும் இருந்தால்தான்
$(\lambd[x][N]M)$ பொருளுடையது. இந்தத் தகவலை
$\lambd$-கலனத்துடன் சேர்த்தால் \emph{வகையிட்ட}
$\lambd$-கலனம் கிடைக்கிறது. இதில் எல்லா
$(\lambd[x][N]M)$ கோவைகளும் அனுமதிக்கப்படுவதில்லை;
$\lambd[x][N]$ உம் $M$ உம் வகைகளின்படி பொருந்தும்
கோவைகள் மட்டுமே அனுமதிக்கப்படுகின்றன. அதாவது,
$\lambd[x][N]$ இன் வகை $A \lif B$ ஆகவும்
$M$ இன் வகை~$A$ ஆகவும் இருக்க வேண்டும்.
$x$ இன் வகை $A$ ஆகவும் $N$ இன் வகை~$B$ ஆகவும்
இருந்தால்தான் $\lambd[x][N]$ இன் வகை $A \to B$
ஆகும். பொதுவாகவும், எல்லா $(M_1M_2)$ கோவைகளும்
அனுமதிக்கப்படுவதில்லை. அத்தகைய $(M_1M_2)$ சொல்லில்,
$M_1$ இன் வகை $A \to B$
ஆகவும் $M_2$ இன் வகை~$A$ ஆகவும் இருந்தால் மட்டுமே
$(M_1M_2)$ அனுமதிக்கப்படும்; அப்போது அதன் வகை~$B$.

% SEGMENT OLP-0686-S03
இவ்வகையிடல் விதிகள் இயல்பான வருவித்தலில் உள்ள
!!{derivation}s உடன் நேரடியாகப் பொருந்துகின்றன:
வகைகளுக்கு !!{formula}s உம், !!{derivation}s க்கு
$\lambd$-சொற்களும் பொருந்துகின்றன. எடுத்துக்காட்டாக,
$!A \lif !B$ இன் !!a{derivation} க்கு
$\lambd[\typeof{x}{!A}][N]$ என்ற சொல் பொருந்தும்;
அது $!A \lif !B$ என்ற சார்பு வகையுடையது.
இயல்பான வருவித்தலின் அனுமான விதிகள், வகையிட்ட
லாம்டா கலனத்தின் வகையிடல் விதிகளுக்குப் பொருந்துகின்றன.
எடுத்துக்காட்டாக,
\begin{prooftree}
  \AxiomC{$\Discharge{!A}{x}$}
  \DeduceC{$!B$}
  \DischargeRule{\Intro\lif}{x}
  \UnaryInfC{$!A \lif !B$}
  \DisplayProof\bottomAlignProof
  \quad\text{இதற்குப் பொருந்துவது}\quad
  \Axiom$x: !A \fCenter N: !B$
  \UnaryInf$ \fCenter \lambd[\typeof{x}{!A}][N] : !A \lif !B$
\end{prooftree}
வலப்புற விதியின் பொருள் இதுதான்: $x$ இன் வகை~$!A$
ஆகவும் $N$ இன் வகை~$!B$ ஆகவும் இருந்தால்,
$\lambd[\typeof{x}{!A}][N]$ இன் வகை~$!A \lif !B$.

% SEGMENT OLP-0686-S04
$\Elim\lif$ விதிக்கு, சொற்களைச் செயல்படுத்துவதற்கான
வகையிடல் விதி பொருந்துகிறது; அதாவது,
\[
  \AxiomC{$!A \lif !B$}
  \AxiomC{$!A$}
  \RightLabel{\Elim\lif}
  \BinaryInfC{$!B$}
  \DisplayProof
  \quad\text{இதற்குப் பொருந்துவது}\quad
  \Axiom$\fCenter M_1: !A \lif !B$
  \Axiom$\fCenter M_2: !A$
  \BinaryInf$ \fCenter (M_1M_2) : !B$
  \DisplayProof
\]
\Intro\lif{} விதியை உடனே \Elim\lif{} விதி
பின்தொடர்ந்தால், அந்த !!{derivation} ஐ எளிமையாக்கலாம்:
\begin{gather*}
  \AxiomC{$\Discharge{!A}{x}$}
  \DeduceC{$!B$}
  \DischargeRule{\Intro\lif}{x}
  \UnaryInfC{$!A \lif !B$}
  \AxiomC{}
  \DeduceC{$!A$}
  \RightLabel{\Elim\lif}
  \BinaryInfC{$!B$}
  \DisplayProof
  \quad\redone\quad
  \AxiomC{}
  \DeduceC{$!A$}
  \DeduceC{$!B$}
  \DisplayProof
\end{gather*}
இதற்கு லாம்டா சொற்களின் $\beta$-சுருக்கல் பொருந்துகிறது:
\[
(\lambd[\typeof{x}{!A}][N]M) \quad\redone\quad
\Subst{N}{M}{x}.
\]

% SEGMENT OLP-0686-S05
$\land$ க்கான விதிகளுக்கும் ``பெருக்கல்'' வகைகளுக்கும்
இடையிலும், $\lor$ க்கான விதிகளுக்கும் ``கூட்டல்''
வகைகளுக்கும் இடையிலும் இதே போன்ற பொருத்தங்கள் உள்ளன.

எளிய வகையிட்ட $\lambd$-கலனச் சொற்களுக்கும் இயல்பான
வருவித்தலின் !!{derivation}s க்கும் இடையிலான இந்தப்
பொருத்தம் ``கரி--ஹாவர்டு பொருத்தம்'',
``நிறுவல்கள் நிரல்களாக'', அல்லது ``கூற்றுகள் வகைகளாக''
என்று அழைக்கப்படுகிறது. !!{formula}s ஆகிய கூற்றுகள்
வகைகளுக்கும், நிறுவல்கள் சொற்களுக்கும் பொருந்துவதோடு,
மற்ற பொருத்தங்களும் பின்வருமாறு:
\begin{center}
  \begin{tabular}{c c c}
    தருக்கம் & நிரல் \\
    \hline
    கூற்று/!!{formula} & வகை \\
    நிறுவல் & சொல் \\
    எடுகோள் & மாறி \\
    விடுவிக்கப்பட்ட எடுகோள் & கட்டுப்பட்ட மாறி \\
    விடுவிக்கப்படாத எடுகோள் & கட்டற்ற மாறி \\
    $!A \lif !B$ & சார்பு வகை \\
    $!A \land !B$ & பெருக்கல் வகை \\
    $!A \lor !B$ & கூட்டல் வகை \\
    $\lfalse$ & வெற்று வகை \\
    \hline
  \end{tabular}
\end{center}

கரி--ஹாவர்டு பொருத்தம் தானியங்கிய நிறுவல்
உதவியாளர்களுக்கும் நிரல்களின் வகைச் சரிபார்ப்பிகளுக்கும்
அடித்தளமான கருத்துகளில் ஒன்று. ஏனெனில், ஒரு கூற்றை
நிறுவும் சான்றைச் சரிபார்ப்பது, அதற்குரிய நிரலான சொல்
அறிவிக்கப்பட்ட வகையுடையதா எனச் சரிபார்ப்பதற்குச் சமம்.
\end{document}
